% !TEX root = ../main.tex
\chapter{Pseudo-exam and second reading}\label{app:exam}
\begin{paracol}{2}

The long exercises of each chapter, moved here verbatim so that a first reading of the chapters is not held
up by them. Use this appendix as a second reading or, once the course is finished, as a cumulative
self-test: retrieval spaced weeks after first study is among the best-supported ways to make material stick
(Cepeda et al.\ 2006; Adesope et al.\ 2017).

% ------------------------------------------------------------------
\section{Logic of propositions and predicates (\cref{ch:logic})}\label{exam:ch01}

\begin{exercise}[NAND is functionally complete]\label{ex:lo-nand-full}
Show that every unary and every binary operator of propositional logic can be written using only the
NAND operator $\uparrow$. Start with $\neg p$, then $p\wedge q$, $p\vee q$, $p\Rightarrow q$ and
$p\Leftrightarrow q$; then argue that every operator of any order $n$ can be built from these.
\end{exercise}

% ------------------------------------------------------------------
\section{Classification of sets (\cref{ch:l03})}\label{exam:ch03}

\begin{exercise}[Well-definedness of the addition on $\mathbb Z$]\label{ex:cl-wd-Z}
On $\N\times\N$ let $(m,n)\sim(p,q):\Leftrightarrow m+q=p+n$. Show that $\sim$ is an equivalence relation (transitivity
included) and that $[(m,n)]+_{\mathbb Z}[(p,q)]:=[(m+p,\,n+q)]$ does not depend on the chosen representatives. Check that
$[(0,0)]$ is neutral and that $[(n,0)]+_{\mathbb Z}[(0,n)]=[(0,0)]$ for every $n\in\N$. Define a multiplication on $\mathbb Z$ by
inheriting the one of $\N$, and show that it is well defined (see \cref{def:cl-Z,def:cl-addZ}).
\end{exercise}

\begin{exercise}[The rational numbers and the fool's addition]\label{ex:cl-wd-Q}
On $\mathbb Z\times\mathbb Z^*$ let $(x,y)\sim(u,v):\Leftrightarrow xv=uy$. Show that $\sim$ is an equivalence relation, and that
$[(x,y)]+_\Q[(u,v)]:=[(xv+yu,\,yv)]$ and $[(x,y)]\cdot_\Q[(u,v)]:=[(xu,\,yv)]$ are well defined (note that $yv\ne0$). Show that
$[(x,y)]\,\tilde+\,[(u,v)]:=[(x+u,\,y+v)]$ is \emph{not} well defined by finding representatives that give different results
(see \cref{def:cl-Q,def:cl-opsQ,rem:cl-fool}). Show that $x\mapsto[(x,1)]$ is injective and compatible with the operations.
\end{exercise}

% ------------------------------------------------------------------
\section{Topological spaces: construction and purpose (\cref{ch:l04})}\label{exam:ch04}

\begin{exercise}[The topology inherited by a quotient set]\label{ex:to-quotient-full}
Let $(M,\mathcal O)$ be a topological space and $\sim$ an equivalence relation on $M$ (see \cref{def:cl-equiv,def:cl-quotient}). Define on the
quotient set $M/{\sim}$ a topology inherited from $\mathcal O$ (the problem-sheet exercise of the lecture does not give the definition: it
asks to find it), and show that it is a topology. Apply it to $M=\R$ with the standard topology and $x\sim y:\Leftrightarrow y-x\in2\pi\mathbb Z$, the
circle as a quotient set (see \cref{ex:to-quotient}).
\end{exercise}

% ------------------------------------------------------------------
\section{Topological spaces: some heavily used invariants (\cref{ch:l05})}\label{exam:ch05}

\begin{exercise}[Limits in Hausdorff spaces and the Zariski topology]\label{ex:iv-limit-full}
Prove that in a Hausdorff space every convergent sequence has a unique limit point (\cref{thm:iv-limit}). Show that the chaotic topology on a set with at least two
points is not $T_1$. Look up the Zariski topology and show that it is $T_1$ but not $T_2$ (\cref{ex:iv-zariski}).
\end{exercise}

\begin{exercise}[{$[0,1]$} is connected]\label{ex:iv-interval-full}
Prove that $[0,1]$ with the topology induced from the standard topology of $\R$ is connected (\cref{thm:iv-interval}); not proved in class, ``a little work''.
\end{exercise}

\begin{exercise}[The fundamental group is a group]\label{ex:iv-group-full}
Show that homotopy is an equivalence relation on curves with fixed end points (transitivity included), that $[\gamma]\bullet[\delta]:=[\gamma*\delta]$ is well defined
(independent of the representatives), associative, has the class of the constant loop as neutral element and $[\gamma\circ(1-\cdot)]$ as inverse
(\cref{def:iv-pi1,thm:iv-group,ex:iv-wd}).
\end{exercise}

% ------------------------------------------------------------------
\section{Topological manifolds and manifold bundles (\cref{ch:l06})}\label{exam:ch06}

\begin{exercise}[The Möbius strip is locally trivial but not trivial]\label{ex:mf-moebius-full}
For the Möbius strip as a bundle over $S^1$ with fibre $[-1,1]$ (\cref{ex:mf-moebius}) find open sets $U\subseteq S^1$ over which the restricted bundle is isomorphic to a product bundle, so that it is locally trivial, and show that the bundle
is not trivial (\cref{def:mf-trivial,ex:mf-cyl-moeb}). Write a section locally as a map $U\to[-1,1]$ (\cref{ex:mf-loctriv}).
\end{exercise}

\begin{exercise}[Pulled-back sections]\label{ex:mf-pullback-full}
Let $(E,\pi,M)$ be a bundle with a section $\sigma$ and $f\colon M'\to M$ a map. Show that $\sigma'(m'):=(m',\sigma(f(m')))$ is a section of the pull-back bundle, that it is continuous if $f$ and $\sigma$ are, and that it is the only
section $\sigma'$ with $u\circ\sigma'=\sigma\circ f$ (\cref{def:mf-pullback,prop:mf-pullback-section,ex:mf-pullback-sec}).
\end{exercise}

\begin{exercise}[Continuity of maps and curves through charts]\label{ex:mf-cont-charts-full}
Let $M,N$ be topological manifolds and $\phi\colon M\to N$ a map. Show that if $\phi$ is continuous, then $y\circ\phi\circ x^{-1}$ is continuous for all charts $(U,x)$ of $M$ and $(V,y)$ of $N$ with $\phi(U)\subseteq V$; and conversely that $\phi$ is continuous if every point of $M$ lies in the domain
of a chart $(U,x)$ for which there is a chart $(V,y)$ of $N$ with $\phi(U)\subseteq V$ and $y\circ\phi\circ x^{-1}$ continuous [?: the precise statement of the problem sheet is not in the lecture]. Deduce the criterion for a curve $\gamma\colon\R\to M$ of \cref{sec:mf-curve}
(\cref{ex:mf-cont-charts}).
\end{exercise}

\flushside
\end{paracol}
